% status: SKELETON
% Chapter 16 of the second edition. Plan: docs/STRUCTURE.md. Rules: AUTHORING.md.
\chapter{Paged KV Memory}\label{ch:paged-kv}

\begin{ChapterIntent}
When every request reserves contiguous KV memory for its maximum length,
most of the reservation is never used and capacity --- and with it batch size
and throughput --- collapses. Paged KV memory borrows virtual memory's answer:
fixed-size blocks, a per-request block table and a pool that allocates on
demand. This chapter builds a block pool with reference counts and
copy-on-write, rewrites attention to read through a block table, prices the
indirection, and discusses block size, fragmentation and eviction.
\end{ChapterIntent}

\OpeningSection{Most of the reservation, never used}

\NapkinSection{Napkin math: fragmentation under real length distributions}

\section{Reservation and fragmentation}

\section{Blocks and block tables}

\section{The block pool and its allocator}

\section{Sharing: reference counts and copy-on-write}

\section{Attention through a block table}

\section{Choosing the block size}

\section{Eviction and swapping at block granularity}

\LedgerSection

\LabSection{a block pool with copy-on-write, checked against a flat cache}

\HappenedSection{the paged path and its release gate}

\FrontierSection{}

\ExercisesSection
